% Standalone section; requires amsmath, amssymb, booktabs, graphicx, tikz, and hyperref.
% For the graph: \usetikzlibrary{arrows.meta,positioning}
\section{Lean correspondence}
\label{app:lean-correspondence}

The accompanying Lean development follows the dependency structure of
Sections~4 and~5.  Figure~\ref{fig:lean-dependencies} gives the dependency
graph (arrows point to checked ingredients), while
Table~\ref{tab:lean-correspondence} gives the statement-by-statement
correspondence.  In the figure, green boxes are paper-facing coefficient
cores and blue boxes are standalone bridge or boundary proofs.  The number in
parentheses is the number of Lean declarations introduced with
\texttt{theorem} or \texttt{lemma}; the disjoint counts sum to 36.

All 36 declarations are checked with Lean~4.32.2 and Mathlib~4.32.2.  The
source contains no \texttt{sorry}, \texttt{admit}, or project-specific
\texttt{axiom}.  A dedicated kernel audit prints the axiom dependencies of the
public declarations.  The word ``core'' in the graph is important: the current
development verifies the finite coefficient and order-theoretic consequences,
not the general-measure discretisation, differentiated R\'enyi integral, or
high-dimensional localisation arguments.  Those analytic boundaries are
listed explicitly in the last row of Table~\ref{tab:lean-correspondence}, and
are not hidden as hypotheses or axioms.

\begin{figure}[t]
  \centering
  \resizebox{\textwidth}{!}{%
  \begin{tikzpicture}[
      x=1cm,y=1cm,
      box/.style={draw,rounded corners,align=center,inner sep=3pt,font=\small,
                  minimum width=2.7cm},
      paper/.style={box,fill=green!8,draw=green!60!black},
      boundary/.style={box,fill=blue!9,draw=blue!65!black},
      every edge/.style={draw,-{Latex[length=1.6mm]}}]
    \node[paper] (l43) at (-9,3.2) {Lemma 4.3 core\\\scriptsize(2 Lean theorems)};
    \node[paper] (p48) at (-6,3.2) {Proposition 4.8 core\\\scriptsize(1 Lean theorem)};
    \node[paper] (p51) at (-3,3.2) {Proposition 5.1 core\\\scriptsize(4 Lean theorems)};
    \node[paper] (p523) at (0,3.2) {Propositions 5.2--5.3 core\\\scriptsize(2 Lean theorems)};
    \node[paper] (p56) at (3,3.2) {Proposition 5.6 core\\\scriptsize(1 Lean theorem)};
    \node[paper] (nt) at (6,3.2) {Negative tropical case, \S5\\\scriptsize(2 Lean theorems)};
    \node[paper] (p57) at (9,3.2) {Proposition 5.7 core\\\scriptsize(2 Lean theorems)};

    \node[boundary] (appa) at (-7.5,0) {Appendix A implication\\\scriptsize(8 Lean theorems)};
    \node[boundary] (b2) at (-2.5,0) {Lemma 4.2 / Lemma B.2\\\scriptsize(9 Lean theorems)};
    \node[boundary] (b15) at (2.5,0) {Lemma B.15\\\scriptsize(2 Lean theorems)};
    \node[boundary] (b6) at (7.5,0) {Appendix B.6\\\scriptsize(3 Lean theorems)};

    \draw (l43) edge[out=-90,in=90] (appa);
    \draw (l43) edge[out=-90,in=90] (b2);
    \draw (p48) edge[out=-90,in=90] (appa);
    \draw (p51) edge[out=-90,in=90] (b2);
    \draw (p51) edge[out=-90,in=90] (b6);
    \draw (p523) edge[out=-90,in=90] (b2);
    \draw (p523) edge[out=-90,in=90] (b6);
    \draw (p56) edge[out=-90,in=90] (b2);
    \draw (p56) edge[out=-90,in=90] (b6);
    \draw (nt) edge[out=-90,in=90] (b2);
    \draw (nt) edge[out=-90,in=90] (b15);
    \draw (nt) edge[out=-90,in=90] (b6);
    \draw (p57) edge[out=-90,in=90] (b2);
    \draw (p57) edge[out=-90,in=90] (b6);
  \end{tikzpicture}%
  }
  \caption{Numbered manuscript dependency graph for the checked Lean core.
  Parenthetical counts total 36 theorem/lemma declarations.}
  \label{fig:lean-dependencies}
\end{figure}

\begin{table}[t]
\centering
\scriptsize
\begin{tabular}{@{}p{0.24\linewidth}p{0.49\linewidth}r p{0.13\linewidth}@{}}
\toprule
Manuscript statement & Lean declaration(s) & Count & Status\\
\midrule
Appendix A, conditioning implication &
\texttt{isSuperadditive\_of\_concave\_posHomogeneous};
\texttt{isSubadditive\_of\_convex\_posHomogeneous};
\texttt{superadditive\_finset\_sum};
\texttt{subadditive\_finset\_sum};
\texttt{columnLift\_pushConditioning\_ge};
\texttt{columnLift\_pushConditioning\_le};
\texttt{conditioning\_monotone\_of\_concave};
\texttt{conditioning\_antitone\_of\_convex} & 8 & Formalized\\
Lemma 4.2 / Lemma B.2, coefficient-curvature consequences &
The nine declarations in \texttt{BoundaryProofs.Curvature} & 9 & Formalized core\\
Lemma 4.3, finite coefficient sign implications &
\texttt{positiveTemperate\_coefficients\_nonnegative};
\texttt{negativeTemperate\_two\_coefficient\_curvature} & 2 & Formalized core\\
Proposition 4.8, coefficient consequence &
\texttt{derivation\_supported\_below\_has\_no\_upper\_obstruction} & 1 & Formalized core\\
Proposition 5.1, coefficient conclusions &
\texttt{positiveTemperate\_shannon\_atom\_zero};
\texttt{positiveTemperate\_no\_upper\_moment};
\texttt{positiveTemperate\_lowerMoment\_le\_one};
\texttt{positiveTemperate\_necessary} & 4 & Formalized core\\
Propositions 5.2--5.3, exceptional moment consequence &
\texttt{negativeTemperate\_truncated\_moment};
\texttt{negativeTemperate\_exceptional\_moment} & 2 & Formalized core\\
Proposition 5.6, two-upper-cell contradiction &
\texttt{negativeTemperate\_atMostOneUpperCell} & 1 & Formalized core\\
Negative tropical necessity in Section 5 &
\texttt{negativeTropical\_moment\_nonnegative};
\texttt{negativeTropical\_exceptional\_moment} & 2 & Formalized core\\
Proposition 5.7, derivation necessity &
\texttt{derivation\_no\_positive\_upper\_tail};
\texttt{derivation\_necessary} & 2 & Formalized core\\
Lemma B.15, exact stationarity &
\texttt{stationarityCorrection\_dot};
\texttt{stationarityCorrection\_tendsto} & 2 & Formalized\\
Appendix B.6, finite midpoint contradictions &
\texttt{not\_concave\_of\_strict\_midpoint\_valley};
\texttt{not\_convex\_of\_strict\_midpoint\_peak};
\texttt{not\_quasiConvex\_of\_strict\_midpoint\_peak} & 3 & Formalized\\
\midrule
Proposition 4.4 and analytic Lemmas B.1, B.3--B.13 &
General-measure approximation, differentiation, and localisation & 0 & Analytic boundary\\
\midrule
\multicolumn{2}{r}{\textbf{Checked theorem/lemma declarations}} & \textbf{36} & \\
\bottomrule
\end{tabular}
\caption{Correspondence between numbered manuscript statements and Lean
declarations.  Counts are disjoint, so the total is the number of checked
theorem/lemma declarations in the project.}
\label{tab:lean-correspondence}
\end{table}
